\input{../style}

\title{{\Large Programmation sur Processeur Graphique -- GPGPU }\\
\vspace{-0.4em}
{\large TD 1 : mon premier programme en CUDA}\\
{\small Centrale Nantes}\\
\small P.-E. Hladik, \small{\texttt{pehladik@ec-nantes.fr}}\\
---\\
{\scriptsize  Version 1.0.0 (\today)}\\
}

\begin{document}

\maketitle

%\begin{table}[htdp]
%\caption{Suivi des modifications}
%\begin{center}
%\begin{tabular}{|c|c|c|p{8cm}|}
%\hline
%Date & Version & Auteur & Modifications\\
%\hline
%JJ/MM/12 & 0.3.X & P.-E. Hladik&  \\
%\hline
%\end{tabular}
%\end{center}
%\label{default}
%\end{table}%

\begin{ad}[]{}
 \begin{itemize}
 	Pour travailler à distance sur la Jetson, voir le guide d'installation sur \href{https://gitlab.univ-nantes.fr/hladik-pe-1/ecn-gpu-cours/-/blob/main/HowTo/Hwt1-remote/GPGPU_Hwt_remote.pdf?ref_type=heads}{gitlab}.
\end{itemize}
\end{ad}


%%%%%%%%%%%%%%%%%%%%%%%%%
  \section{Parallèle Hello World}
%%%%%%%%%%%%%%%%%%%%%%%%%

\begin{objectif}[]{}
 \begin{itemize}
	\item Compiler et exécuter un code en CUDA.
\end{itemize}
\end{objectif}

\begin{wtodo}[Hello world!]{}
\begin{enumerate}
\item Récupérer le fichier {\tt hello.cu} sur le dépôt et déposer le sur la Jetson (un simple glisser/déposer marche très bien avec VSC Remote).
\item Compiler le fichier {\tt hello.cu}.
\item Lancer l'exécutable, vous venez d'exécuter votre premier code sur un GPU !
\end{enumerate}

\end{wtodo}

\begin{howto}[compiler un code CUDA]{}
Pour compiler un fichier en CUDA (extension communément utilisée {\tt .cu}) il vous faut un compilateur dédié, par exemple {\tt nvcc}. Vous pouvez consulter la documentation sur :

 \url{https://docs.nvidia.com/cuda/cuda-compiler-driver-nvcc/index.html}
 
\noindent La commande de base pour compiler est :
 
 {\tt nvcc -o outputFile inputFile.cu}

avec {\tt outputFile} le nom du fichier de sortie et {\tt inputFile} le fichier source à compiler.
\end{howto} 

\begin{howto}[lancer l'exécutable]{}
L'exécution d'un code CUDA se fait exactement de la même manière que pour un programme standard en faisant un simple appel à {\tt ./outputFile}.
\end{howto} 

%%%%%%%%%%%%%%%%%%%%%%%%%
  \section{Connaître sa machine}
%%%%%%%%%%%%%%%%%%%%%%%%%

\begin{objectif}[]{}
 \begin{itemize}
	\item Afficher les caractéristiques du GPU
\end{itemize}
\end{objectif}

\begin{wtodo}[Lecture des paramètres du GPU]{}
\begin{enumerate}
\item Récupérer le template {\tt properties.cu}, compilez exécutez le. Que signifie le résultat affiché (faite une recherche sur internet) ?
\item Modifier le code pour connaître le nombre de GPU sur la machine et pour afficher le nom de la carte GPU que vous utilisez. Compiler, exécuter.
\item Ajouter le code nécessaire pour afficher le nombre maximum de threads pouvant résider sur un multiprocesseur ({\tt maxThreadsPerMultiProcessor})
\end{enumerate}

\end{wtodo}

\begin{howto}[La structure {\tt struct cudaDeviceProp }]{}
La structure {\tt cudaDeviceProp} contient un ensemble de champs\footnote{\url{https://docs.nvidia.com/cuda/cuda-runtime-api/structcudaDeviceProp.html}} décrivant un {\it device}. La structure est récupérée à l'aide de la fonction {\tt cudaGetDeviceProperties}\footnote{\url{https://docs.nvidia.com/cuda/cuda-runtime-api/group__CUDART__DEVICE.html}}.
\end{howto} 


\begin{howto}[Affichage des champs de la structure {\tt cudaDeviceProp}]{}
Pour afficher les paramètres du GPU\footnote{\url{https://docs.nvidia.com/cuda/cuda-c-programming-guide/index.html\#compute-capabilities}} il suffit de lire le champ de la structure {\tt cudaDeviceProp} comme n'importe quelle structure en C, par exemple :\vspace{-0.5em}
\begin{Verbatim}
int deviceCount;
cudaGetDeviceCount(&deviceCount);
int device;
for (device = 0; device < deviceCount; ++device) {
    cudaDeviceProp deviceProp;
    cudaGetDeviceProperties(&deviceProp, device);
    printf("Device %d has compute capability %d.%d.\n",
           device, deviceProp.major, deviceProp.minor);
}
\end{Verbatim}
\end{howto} 
\end{document}